\subsection{L(X) 的齐次分量的维数}

设 X 是一个集合， \(\alpha\) 是 \(\mathbf{N}^{(X)}\) 的一个元素，d 是一个 >0 的整数。我们记 \(d|\alpha\) ，如果存在 \(\beta\in\mathbf{N}^{(X)}\) 使得 \(\alpha=d\beta\) 。这时唯一的元素 \(\beta\) 记作 \(\alpha/d\) 。

引理 1. 设 \(n\) 是一个整数 \(>0\) ， \(T_1, \ldots, T_n\) 是不定元， \(u_1, \ldots, u_n\) 是 \(\mathbf{Z}\) 的元素。设 \((c(\alpha))_{\alpha \in \mathbf{N}^n - \{0\}}\) 是 \(\mathbf{Z}\) 的元素的一个族，使得

\[
1 - \sum_ {i = 1} ^ {n} u _ {i} T _ {i} = \prod_ {\alpha \neq 0} (1 - T ^ {\alpha}) ^ {c (\alpha)}.\tag{7}
\]

对一切 \(\alpha \in \mathbf{N}^n - \{0\}\) ，

\[
c (\alpha) = \frac {1}{| \alpha |} \sum_ {d | \alpha} \mu (d) \frac {(| \alpha | / d) !}{(\alpha / d) !} u ^ {\alpha / d}\tag{8}
\]

其中 \(\mu\) 是 Möbius 函数（附录）。

在两边取对数后 (《代数》， 第 IV 章， § 6， no. 9)，公式 (7) 等价于：

\[
\log \left(1 - \sum_ {i = 1} ^ {n} u _ {i} T _ {i}\right) = \sum_ {\alpha \neq 0} c (\alpha) \log (1 - T ^ {\alpha}).\tag{9}
\]

而

\[
\begin{array}{r l} - \log \left(1 - \sum_ {i = 1} ^ {n} u _ {i} \mathrm{T} _ {i}\right) & = \sum_ {j \geqslant 1} \frac {1}{j} \left(\sum_ {i = 1} ^ {n} u _ {i} \mathrm{T} _ {i}\right) ^ {j} \\ & = \sum_ {j \geqslant 1} \frac {1}{j} \sum_ {| \beta | = j} \frac {| \beta | !}{\beta !} u ^ {\beta} \mathrm{T} ^ {\beta} \\ & = \sum_ {| \beta | > 0} \frac {1}{| \beta |} \frac {| \beta | !}{\beta !} u ^ {\beta} \mathrm{T} ^ {\beta}. \end{array}\tag{10}
\]

另一方面

\[
\begin{array}{r l} - \sum_ {\alpha \neq 0} c (\alpha) \log (1 - T ^ {\alpha}) & = \sum_ {| \alpha | > 0, k \geqslant 1} \frac {1}{k} c (\alpha) T ^ {k \alpha} \\ & = \sum_ {| \beta | > 0, k | \beta} \frac {1}{k} c \left(\frac {\beta}{k}\right) T ^ {\beta}. \end{array}\tag{11}
\]

因此 (7) 等价于

\[
\sum_ {k \mid \beta} \left| \frac {\beta}{k} \right| c \left(\frac {\beta}{k}\right) = \frac {| \beta | !}{\beta !} u ^ {\beta} \quad \text { for   all } \beta \in \mathbf {N} ^ {n} - \{0 \}.\tag{12}
\]

设 \(\Lambda\) 是由 \((\lambda_1, \lambda_2, \ldots, \lambda_n) \in \mathbf{N}^n - \{0\}\) 组成的集合，其中 \(\lambda_1, \lambda_2, \ldots, \lambda_n\) 的最大公因子等于 1。 \(\mathbf{N}^n - \{0\}\) 的每个元素都可以唯一地写成

形式 \(m\lambda\) ，其中 \(m\) 是一个整数 \(\geqslant 1\) 且 \(\lambda \in \Lambda\) 。条件 (12) 等价于

\[
\sum_ {k \mid m} \left| \frac {m \lambda}{k} \right| c \left(\frac {m \lambda}{k}\right) = \frac {(m | \lambda |) !}{(m \lambda) !} u ^ {m \lambda} \quad \text { for   all } \lambda \in \Lambda \text { and   all } m \geqslant 1.\tag{13}
\]

由 Möbius 反演公式（附录），条件 (13) 等价于

\[
\left| m \lambda \right| c (m \lambda) = \sum_ {d \mid m} \mu (d) \frac {\left| \frac {m \lambda}{d} \right| !}{\left(\frac {m \lambda}{d}\right) !} u ^ {\frac {m \lambda}{d}}\tag{14}
\]

其中一切 \(\lambda\in\Lambda\) 与一切 \(m\geqslant1\) 。

定理 2. 设 X 是一个有限集，并且 \(n = \operatorname{Card}(X)\) 。

\begin{enumerate}
\def\labelenumi{(\alph{enumi})}
\tightlist
\item
  对每个整数 \(r \geqslant 1\) ，K-模 \(\mathbf{L}^r(\mathbf{X})\) 都是自由的，其秩为
\end{enumerate}

\[
c (r) = \frac {1}{r} \sum_ {d | r} \mu (d) n ^ {r / d},\tag{15}
\]

其中 μ 是 Möbius 函数。

\begin{enumerate}
\def\labelenumi{(\alph{enumi})}
\setcounter{enumi}{1}
\tightlist
\item
  对一切 \(\alpha \in \mathbf{N}^{(\mathbf{X})} - \{0\}\) ，K-模 \(\mathbf{L}^{\alpha}(\mathbf{X})\) ( \(\S 2\) ， no. 6) 都是自由的，其秩为
\end{enumerate}

\[
c (\alpha) = \frac {1}{| \alpha |} \sum_ {d | \alpha} \mu (d) \frac {(| \alpha | / d) !}{(\alpha / d) !}.\tag{16}
\]

我们已经知道，模 \(L^{r}(\mathbf{X})\) （其中 \(r \in \mathbf{N}\) ）与 \(L^{\alpha}(\mathbf{X})\) （其中 \(\alpha \in \mathbf{N}^{\mathbf{X}}\) ）都是自由的 (§ 2， no. 11， 定理 1 的推论)。考察 A(X) 的多重分次 \((A^{\alpha}(X))_{\alpha \in \mathbf{N}^{\mathbf{X}}}\) ，它由典范同态 \(\phi\) （从 Mo(X) 到 \(\mathbf{N}^{\mathbf{X}}\) 中）所定义 (《代数》， 第 III 章， § 3， no. 1， 例 3)；于是由 no. 1 的注 3 有 \(A^{\alpha}(\mathbf{X}) \cap L(\mathbf{X}) = L^{\alpha}(\mathbf{X})\) 。对 \(\alpha \in \mathbf{N}^{\mathbf{X}}\) ，K-模 \(A^{\alpha}(\mathbf{X})\) 以这样的字的集合为基：X 的每个字母 \(x\) 在其中出现 \(\alpha(x)\) 次。设 \(d(\alpha)\) 是这些字的个数，即 \(A^{\alpha}(\mathbf{X})\) 的秩；我们将以两种不同的方式计算形式幂级数

它由

\[
\mathbf {P} ((\mathrm{T} _ {x}) _ {x \in \mathbf {X}}) \in \mathbf {Z} [ [ (\mathrm{T} _ {x}) _ {x \in \mathbf {X}} ] ]\tag{17}
\]

\[
\mathrm{P} (\mathrm{T}) = \sum_ {\alpha \in \mathbf {N} ^ {\mathbf {X}}} d (\alpha) \mathrm{T} ^ {\alpha}.
\]

\begin{enumerate}
\def\labelenumi{(\arabic{enumi})}
\tightlist
\item
  我们有
\end{enumerate}

\[
\mathrm{P} (\mathrm{T}) = \sum_ {m \in \mathrm{Mo} (\mathrm{X})} \mathrm{T} ^ {\phi (m)} = \sum_ {r = 0} ^ {\infty} \sum_ {x _ {1}, \dots , x _ {r}} \mathrm{T} _ {x _ {1}} \dots \mathrm{T} _ {x _ {r}} = \sum_ {r = 0} ^ {\infty} \left(\sum_ {x \in \mathrm{X}} \mathrm{T} _ {x}\right) ^ {r}
\]

从而

\[
\mathrm{P} (\mathrm{T}) = \left(1 - \sum_ {x \in \mathbf {X}} \mathrm{T} _ {x}\right) ^ {- 1}.\tag{18}
\]

\begin{enumerate}
\def\labelenumi{(\arabic{enumi})}
\setcounter{enumi}{1}
\tightlist
\item
  对一切 \(\alpha \in \mathbf{N}^{\mathbf{X}} - \{0\}\) ，设 \((e_{\alpha,j})_{1 \leqslant j \leqslant c(\alpha)}\) 是 \(\mathrm{L}^{\alpha}(\mathbf{X})\) 的一个基，并给集合 I 一个全序，I 由有序对 \((\alpha, j)\) （满足 \(\alpha \in \mathbf{N}^{\mathbf{X}} - \{0\}\) 且 \(1 \leqslant j \leqslant c(\alpha)\) ）组成。由 no. 1 的定理 1 与 Poincaré-Birkhoff-Witt 定理 (第 I 章， § 2， no. 7， 定理 1 的推论 3)，诸元素
\end{enumerate}

\[
y _ {m} = \prod_ {(\alpha , j) \in \mathbb {I}} (e _ {\alpha , j}) ^ {m (\alpha , j)},
\]

（其中指标 \(\pmb{m}\) 取遍 \(\mathbf{N}^{(I)}\) ）构成 \(A(\mathbf{X})\) 的一个基。每个 \(y_{m}\) 都有一个多重次数 \(\sum_{(\alpha, j) \in I} \pmb{m}(\alpha, j)\alpha\) 。用 \(u(\pmb{m})\) 表示这个多重次数。由此得出

\[
\begin{array}{l} \mathrm{P(T)} = \sum_ {m \in \mathrm{N} ^ {(I)}} \mathrm{T} ^ {u (m)} = \sum_ {m \in \mathrm{N} ^ {(I)}} \prod_ {(\alpha , j) \in I} \mathrm{T} ^ {m (\alpha , j) \alpha} \\ = \prod_ {(\alpha , j) \in I} \sum_ {r = 0} ^ {\infty} \mathrm{T} ^ {r \alpha} = \prod_ {(\alpha , j) \in I} (1 - \mathrm{T} ^ {\alpha}) ^ {- 1}, \end{array}
\]

从而最后

\[
\mathrm{P} (\mathrm{T}) = \prod_ {\alpha \in \mathbf {N} ^ {\mathbf {X}} - \{0 \}} (1 - \mathrm{T} ^ {\alpha}) ^ {- c (\alpha)}.\tag{19}
\]

比较 (18) 与 (19)，我们得到

\[
1 - \sum_ {x \in \mathbf {X}} T _ {x} = \prod_ {\alpha \in \mathbf {N} ^ {\mathbf {X}} - \{0\}} (1 - T ^ {\alpha}) ^ {c (\alpha)}.\tag{20}
\]

于是引理 1 给出 (b)。

若现在在公式 (20) 中把同一个不定元 U 代入诸 \(T_{x}\) （对 \(x \in X\) ），则得到

\[
1 - n \mathrm{U} = \prod_ {\alpha \in \mathbf {N} ^ {\mathrm{X}} - \{0\}} (1 - \mathrm{U} ^ {| \alpha |}) ^ {c (\alpha)} = \prod_ {r > 0} (1 - \mathrm{U} ^ {r}) ^ {c (r)}.
\]

再次应用引理 1，我们推出 (a)。

例. 我们有

\[
\begin{array}{l l} c (1) = n, & c (2) = \frac {1}{2} (n ^ {2} - n), \qquad c (3) = \frac {1}{3} (n ^ {3} - n), \\ c (4) = \frac {1}{4} (n ^ {4} - n ^ {2}), & c (5) = \frac {1}{5} (n ^ {5} - n), \qquad c (6) = \frac {1}{6} (n ^ {6} - n ^ {3} - n ^ {2} + n). \end{array}
\]

注. 设 X 是一个集合并设 \(\alpha \in \mathbf{N}^{(X)}\) ；自由 K-模 \(\mathbf{L}^{\alpha}(X)\) 的秩也由公式 (16) 给出。这由定理 2 (b) 以及 § 2 no. 6 的命题 4 立即得出。

